\clearpage
\VerdierPhysicalPageAnchor{0155}%
\begingroup\small
\setlength{\parskip}{0pt}\setlength{\emergencystretch}{1.5em}

\VerdierSemanticAnchor{ch3-3-2-2-admissible}
$\mathcal D^\wedge$ の対象が二つの表現可能関手の間の射の像であるとき、
これを\emph{許容対象}という。関手
$\operatorname{Im}\colon\mathcal{F}\ell_1(\mathcal D)\to\mathcal D^\wedge$
を
\VerdierSemanticAnchor{ch3-3-2-2-1}
\[
\operatorname{Im}\bigl(X\xrightarrow fY\bigr)=h(f)\text{ の像}.
\tag{3.2.2.1}
\]
で定める。$\operatorname{Im}$ は加法的で、無視可能射を零射へ送るので、
一意的に加法関手（同じ記号を濫用する）
\VerdierSemanticAnchor{ch3-3-2-2-2}
\[
\operatorname{Im}\colon\mathrm A(\mathcal D)\longrightarrow\mathcal D^\wedge.
\tag{3.2.2.2}
\]
を経由する。合成関手
$\operatorname{Im}\circ\mathcal H\colon
\mathcal D\longrightarrow\mathrm A(\mathcal D)
\longrightarrow\mathcal D^\wedge$
はほかならぬ関手 $h$ であり、
$\mathcal D^\wedge$ の対象が $\mathrm A(\mathcal D)$ の対象の
$\operatorname{Im}$ による像と同型である
ことと、それが許容対象であることは同値である。

\VerdierSemanticAnchor{ch3-3-2-3}
\medskip\noindent\textbf{命題 3.2.3.}\quad
\textit{関手 $\operatorname{Im}\colon\mathrm A(\mathcal D)\to
\mathcal D^\wedge$ は充満忠実である。}

\VerdierSemanticAnchor{ch3-3-2-3-morphism-square}
\[
\begin{tikzcd}[column sep=4em,row sep=2.4em]
X\arrow[r,"f"]\arrow[d,"g"']&Y\arrow[d,"g'"]\\
X'\arrow[r,"f'"']&Y'
\end{tikzcd}
\qquad=(g;g')
\]
を $\mathcal{F}\ell_1(\mathcal D)$ の射とし、その像を
$\operatorname{Im}((g;g'))$ と書く。図式
\VerdierSemanticAnchor{ch3-3-2-3-canonical-diagram}
\[
\begin{tikzcd}[column sep=4em,row sep=2.2em]
X\arrow[r,"\operatorname{id}_X"]\arrow[d,"\operatorname{id}_X"']&X\arrow[d,"f"]\\
X\arrow[r,"f"]\arrow[d,"f"']&Y\arrow[d,"\operatorname{id}_Y"]\\
Y\arrow[r,"\operatorname{id}_Y"']&Y
\end{tikzcd}
\]
に $\operatorname{Im}$ を適用すると
\[
h(X)\xrightarrow{\operatorname{Im}((\operatorname{id}_X;f))}
\operatorname{Im}(f)\xrightarrow{\operatorname{Im}((f;\operatorname{id}_Y))}h(Y).
\]
を得る。
\VerdierSemanticAnchor{ch3-3-2-3-1}
したがって
\[
\operatorname{Im}((f;\operatorname{id}_Y))\circ
\operatorname{Im}((\operatorname{id}_X;f))=h(f).
\tag{3.2.3.1}
\]
\endgroup
